\section{Deployment and MLOps Practices}

\subsection{Deployment Architecture}

The production system follows a microservices architecture:

\begin{figure}[h]
\centering
\includegraphics[width=\textwidth]{figures/architecture_diagram.png}
\caption{System Architecture: Containerized services with Docker Compose}
\label{fig:architecture}
\end{figure}

\subsubsection{Components}

\begin{enumerate}
    \item \textbf{FastAPI Application}: REST API for predictions
    \begin{itemize}
        \item Port: 8000
        \item Container: bus-classifier-api
        \item Health checks every 30 seconds
    \end{itemize}
    
    \item \textbf{MLflow Tracking Server}: Experiment tracking and model registry
    \begin{itemize}
        \item Port: 5000
        \item Container: bus-classifier-mlflow
        \item Persistent storage via volume mounts
    \end{itemize}
    
    \item \textbf{Streamlit Dashboard}: Interactive visualization and monitoring
    \begin{itemize}
        \item Port: 8501
        \item Standalone or containerized deployment
        \item Connects to API and MLflow
    \end{itemize}
\end{enumerate}

\subsection{Docker Compose Configuration}

Multi-container orchestration with Docker Compose:

\begin{lstlisting}[language=yaml, caption=docker-compose.yml]
version: '3.8'

services:
  api:
    build: .
    container_name: bus-classifier-api
    ports:
      - "8000:8000"
    volumes:
      - ./mlruns:/app/mlruns
      - ./models:/app/models
    environment:
      - MLFLOW_TRACKING_URI=mlruns
      - USE_LOCAL_MODEL=true
    restart: unless-stopped
    healthcheck:
      test: ["CMD", "curl", "-f", "http://localhost:8000/health"]
      interval: 30s
      timeout: 10s
      retries: 3

  mlflow:
    image: python:3.11-slim
    container_name: bus-classifier-mlflow
    ports:
      - "5000:5000"
    volumes:
      - ./mlruns:/mlruns
    command: >
      sh -c "pip install mlflow &&
             mlflow ui --host 0.0.0.0 --port 5000 
                       --backend-store-uri /mlruns"
    restart: unless-stopped
\end{lstlisting}

\subsection{CI/CD Pipeline}

\subsubsection{Continuous Integration}

Automated testing on every commit:

\begin{enumerate}
    \item \textbf{Unit Tests}: pytest runs 19 test cases covering transformers, pipeline, and utilities
    \item \textbf{Integration Tests}: End-to-end pipeline validation
    \item \textbf{API Tests}: FastAPI TestClient validates all endpoints
    \item \textbf{Code Coverage}: pytest-cov ensures >80\% coverage
\end{enumerate}

\subsubsection{Code Quality}

Automated linting and formatting:

\begin{itemize}
    \item \textbf{flake8}: PEP 8 style compliance
    \item \textbf{black}: Automated code formatting
    \item \textbf{isort}: Import statement organization
    \item \textbf{mypy}: Static type checking
\end{itemize}

\subsubsection{Continuous Deployment}

GitHub Actions workflow stages:

\begin{enumerate}
    \item \textbf{Build Stage}: Compile Docker images, check for errors
    \item \textbf{Test Stage}: Run full test suite in container
    \item \textbf{Training Stage}: Retrain model with latest data, log to MLflow
    \item \textbf{Registry Stage}: Register new model version if performance improves
    \item \textbf{Deploy Stage}: Update production containers (manual approval)
\end{enumerate}

\subsection{Model Versioning and Registry}

MLflow Model Registry manages model lifecycle:

\begin{table}[h]
\centering
\caption{Model Registry States}
\begin{tabular}{lp{8cm}}
\toprule
\textbf{Stage} & \textbf{Description} \\
\midrule
None & Newly registered, not yet validated \\
Staging & Under evaluation, not production-ready \\
\textbf{Production} & \textbf{Currently serving predictions (v2)} \\
Archived & Deprecated, kept for historical reference \\
\bottomrule
\end{tabular}
\label{tab:model_stages}
\end{table}

Model transition workflow:
\begin{enumerate}
    \item Train new model → Register as v3 (Stage: None)
    \item Validate on holdout set → Promote to Staging
    \item A/B test against Production model
    \item If superior → Promote to Production, archive old version
\end{enumerate}

\subsection{Data Versioning with DVC}

DVC ensures reproducible data pipelines:

\begin{lstlisting}[language=bash, caption=DVC Pipeline]
# Define pipeline stages
dvc run -n preprocess \
    -d data/raw/passengers.csv \
    -d data/raw/bus.csv \
    -o data/processed/features.csv \
    python src/preprocess.py

dvc run -n train \
    -d data/processed/features.csv \
    -o models/pipeline.joblib \
    python src/train_mlflow.py

# Reproduce pipeline
dvc repro
\end{lstlisting}

Benefits:
\begin{itemize}
    \item Automatic dependency tracking
    \item Cacheable pipeline stages
    \item Shareable across team via Git
    \item Remote storage for large artifacts
\end{itemize}

\subsection{Monitoring and Observability}

\subsubsection{System Monitoring}

\begin{itemize}
    \item \textbf{Health checks}: Docker healthcheck ensures API availability
    \item \textbf{Logging}: Structured logs via Python logging module
    \item \textbf{Metrics}: Prometheus-compatible metrics endpoint (planned)
\end{itemize}

\subsubsection{Model Monitoring}

Dashboard tracks key model health indicators:

\begin{enumerate}
    \item \textbf{Prediction Distribution}: Monitor class balance over time
    \item \textbf{Feature Drift}: Track feature statistics (mean, std, outliers)
    \item \textbf{Performance Metrics}: Live F1, accuracy on recent predictions
    \item \textbf{Latency}: API response time percentiles
\end{enumerate}

\subsubsection{Alerting (Future Work)}

Planned alerting system:
\begin{itemize}
    \item Performance degradation (F1 drops >5\%)
    \item Data drift detection (Kolmogorov-Smirnov test)
    \item API errors (5xx responses >1\%)
    \item Model staleness (>30 days since retraining)
\end{itemize}

\subsection{Security Considerations}

\subsubsection{API Security}

\begin{itemize}
    \item \textbf{Input validation}: Pydantic models prevent malformed requests
    \item \textbf{Rate limiting}: Prevent abuse (planned: 100 req/min per IP)
    \item \textbf{Authentication}: API key-based auth (planned)
    \item \textbf{HTTPS}: TLS encryption for production deployment
\end{itemize}

\subsubsection{Data Privacy}

\begin{itemize}
    \item GPS coordinates never stored persistently
    \item User IDs hashed before logging
    \item GDPR compliance: No PII retention
    \item Anonymized data for model training
\end{itemize}

\subsection{Scalability}

\subsubsection{Horizontal Scaling}

API can scale horizontally behind a load balancer:

\begin{lstlisting}[language=yaml, caption=Kubernetes Deployment (example)]
apiVersion: apps/v1
kind: Deployment
metadata:
  name: bus-classifier-api
spec:
  replicas: 3  # Multiple instances
  selector:
    matchLabels:
      app: bus-classifier
  template:
    spec:
      containers:
      - name: api
        image: bus-classifier-api:latest
        ports:
        - containerPort: 8000
\end{lstlisting}

\subsubsection{Performance Optimization}

\begin{itemize}
    \item Model loaded once at startup (not per request)
    \item Batch prediction support for efficiency
    \item Feature computation cached where possible
    \item Database connection pooling (if database added)
\end{itemize}

\subsection{Cost Analysis}

Estimated monthly operational costs for AWS deployment:

\begin{table}[h]
\centering
\caption{Deployment Cost Estimate (AWS, us-east-1)}
\begin{tabular}{lcc}
\toprule
\textbf{Service} & \textbf{Configuration} & \textbf{Monthly Cost} \\
\midrule
EC2 (t3.medium) & API + MLflow & \$30 \\
ECS Fargate & 2 containers, 1 vCPU & \$45 \\
S3 Storage & 100 GB (data + models) & \$2.30 \\
Data Transfer & 500 GB/month & \$45 \\
\midrule
\textbf{Total} & & \textbf{\$122/month} \\
\bottomrule
\end{tabular}
\label{tab:costs}
\end{table}

For educational/small-scale deployment, free tiers or local hosting suffice.

\subsection{Maintenance and Updates}

\subsubsection{Regular Maintenance}

\begin{itemize}
    \item \textbf{Weekly}: Review MLflow experiments, compare model versions
    \item \textbf{Monthly}: Retrain model with accumulated data
    \item \textbf{Quarterly}: Dependency updates, security patches
    \item \textbf{Annually}: Full system audit, architecture review
\end{itemize}

\subsubsection{Update Strategy}

\begin{enumerate}
    \item Pull latest code from Git
    \item Run full test suite
    \item Build new Docker image
    \item Deploy to staging environment
    \item Smoke tests on staging
    \item Blue-green deployment to production
    \item Monitor for 24 hours
    \item Rollback if issues detected
\end{enumerate}

\subsection{Documentation}

Comprehensive documentation in \texttt{docs/} folder:

\begin{itemize}
    \item \textbf{docs/setup/}: Installation and configuration guides
    \item \textbf{docs/usage/}: User guides and tutorials
    \item \textbf{docs/development/}: Architecture and MLOps roadmap
    \item \textbf{docs/archive/}: Historical phase summaries
    \item \textbf{requirements/README.md}: Modular installation guide
\end{itemize}

All code includes docstrings following Google style guide. API documentation auto-generated via FastAPI's OpenAPI integration.
